% !TeX root = report.tex
\newcommand*{\PathToAssets}{../assets}%
\newcommand*{\PathToOutput}{../output}%

\documentclass[12pt]{article}
\usepackage{my_article_header}
\usepackage{my_common_header}

\begin{document}

\title{
Replication of ``Holding Period Effects in Dividend Strip Returns''
\\{\color{blue} \large Replication Report}
}

\author{
Jie Lin and Zimeng Yi\footnote{University of Chicago}
}

\begin{titlepage}
\maketitle

\doublespacing
\begin{abstract}
We replicate Figure 1 from \cite{golez2024holding}, which plots
12-month maturity interest rates from three sources:
option-implied rates derived from S\&P 500 index options,
zero curve rates from OptionMetrics, and constant maturity Treasury rates from FRED.
Following the methodology in the paper, we estimate option-implied interest rates
using the outer product approach for 1996--2003 and the regression approach for 2004--2022.
Our replication closely matches the paper's findings:
the average difference between the option-implied rate and the zero curve rate is
8.15 basis points in our sample, compared with 7 basis points reported in the paper.
We then extend the sample through 2024 and find that the option-implied rate
falls below the zero curve rate on average in the post-2022 period.
This suggests that the relationship documented in the original sample may have weakened
or reversed in recent years.
\end{abstract}

\end{titlepage}

\doublespacing

\section{Introduction}

\cite{golez2024holding} estimate short-term dividend strip prices
from 27 years of S\&P 500 index options data covering 1996--2022.
A key methodological contribution of the paper is the use of
option-implied interest rates rather than exogenous zero curve rates
when estimating strip prices.
The authors show that even a small error of 6 basis points in the interest rate
can bias half-annual dividend strip returns by 0.86\%.

Figure 1 of the paper plots three 12-month maturity interest rate series:
the option-implied rate, the zero curve rate from OptionMetrics,
and the constant maturity Treasury rate from FRED.
The paper reports that the option-implied rate is on average 7 basis points
higher than the zero curve rate over the full sample,
with the relative difference increasing from 0.64\% in the first half of the sample
to 11.80\% in the second half.

In this replication, we reproduce Figure 1 using WRDS-accessed OptionMetrics
and CRSP data together with Treasury rate data from FRED,
and then extend the sample through 2024
to examine whether the relationship between the two rates
has changed in recent years.

\section{Data}

We use data from the following sources.

\textbf{SPX options.}
We use European S\&P 500 index options from OptionMetrics via WRDS,
covering January 1996 through December 2024.
Following the paper, we use options observed on the last business day of each month,
with maturity of at least 90 days and bid and ask quotes both at or above \$3.
We further restrict the sample to options with moneyness between 0.5 and 1.5.
One important difference from the paper is that the original study uses intraday CBOE data
between 10 a.m.\ and 2 p.m., whereas our replication uses end-of-day OptionMetrics data.
This difference likely accounts for part of the remaining gap between our estimates
and the published values.

\textbf{Zero curve rates.}
We obtain daily zero curve rates from OptionMetrics.
These rates are expressed in continuously compounded form and are linearly interpolated
to obtain a 1-year constant-maturity zero curve rate.

\textbf{Treasury rates.}
We obtain the 1-year constant maturity Treasury rate from FRED
(series DGS1) and convert it to continuously compounded form.

\section{Methodology}

Following Section 2.1 of the paper, we estimate option-implied interest rates
using two methods.

\textbf{Outer Product Approach (1996--2003).}
For each pair of strike prices $(X_i, X_j)$ where $X_i > X_j$, we compute
\begin{equation}
r_f
=
-\frac{1}{\tau}
\ln\left[
\frac{(P(X_i)-C(X_i))-(P(X_j)-C(X_j))}{X_i-X_j}
\right].
\end{equation}
We then take the median across all valid strike pairs for a given date and maturity.

\textbf{Regression Approach (2004--2022).}
For later years, we estimate
\begin{equation}
S_t - C(X) + P(X) = P^{\tau} + \beta X + \epsilon,
\end{equation}
where $\beta = e^{-r_f \tau}$.
The implied interest rate is then recovered as
\begin{equation}
r_f = -\frac{1}{\tau}\ln \hat{\beta}.
\end{equation}

For both the option-implied rate and the zero curve rate,
we linearly interpolate across maturities to obtain a 1-year constant-maturity series.

\section{Results}

\subsection{Figure 1 Replication}

Figure \ref{fig:figure1} replicates Figure 1 from \cite{golez2024holding}.
The three interest rate series track one another closely throughout the sample period.
Consistent with the paper, the option-implied rate is generally slightly above
the zero curve rate.

Our replication produces an average difference of 8.15 basis points
between the option-implied rate and the zero curve rate,
compared with 7 basis points reported in the paper.
The full-sample relative difference in our replication is 3.14\%,
compared with 2.82\% in the paper.
Table \ref{table:figure1_summary} reports the full set of comparison statistics.

\begin{figure}[h]
\centering
\caption{12-Month Interest Rates (1996--2022)}
\includegraphics[width=0.95\linewidth]{\PathToOutput/figure1/figure1.png}
\caption*{
This figure plots three 12-month maturity interest rate series from January 1996
through December 2022: the option-implied rate estimated from put-call parity pairs
of S\&P 500 index options, the zero curve rate from OptionMetrics,
and the 1-year constant maturity Treasury rate from FRED.
All rates are expressed in continuously compounded percentage terms.
}
\label{fig:figure1}
\end{figure}

\begin{table}[h]
\centering
\caption{Figure 1 Summary Statistics}
\begin{tabular}{lrr}
\toprule
Statistic & Our Value & Paper Value \\
\midrule
Average difference (bp) & 8.15 & 7.00 \\
Full-sample relative difference (\%) & 3.14 & 2.82 \\
First-half relative difference (\%) & 0.94 & 0.64 \\
Second-half relative difference (\%) & 12.44 & 11.80 \\
\bottomrule
\end{tabular}

\caption*{
This table compares our replicated summary statistics with those reported in
\cite{golez2024holding}.
The average difference is measured in basis points (bp).
Relative differences are computed relative to the average zero curve rate.
}
\label{table:figure1_summary}
\end{table}

\subsection{Extension: Updated Sample Through 2024}

Figure \ref{fig:figure1_extension} extends the sample through 2024.
The red dashed vertical line marks the end of the original paper's sample period
in December 2022.

In the extended period, the option-implied rate falls below the zero curve rate on average.
Our summary statistics show that the average implied-minus-zero spread is negative
over 2023--2024.
This pattern differs from the original sample and suggests that the relationship
documented in the paper may have weakened or reversed in recent years.

A robustness check indicates that this result is not driven solely by a single outlier month.
Even after removing the most negative month in the extended period,
the average spread remains negative.

\begin{figure}[h]
\centering
\caption{12-Month Interest Rates (1996--2024, Extended Sample)}
\includegraphics[width=0.95\linewidth]{\PathToOutput/figure1_extension/figure1_extension.png}
\caption*{
This figure extends the Figure 1 replication through 2024.
The red dashed vertical line marks the end of the original paper's sample in December 2022.
In the extended period, the option-implied rate is on average below the zero curve rate,
in contrast to the pattern observed during 1996--2022.
}
\label{fig:figure1_extension}
\end{figure}

\begin{table}[h]
\centering
\caption{Extension Summary Statistics}
\begin{tabular}{lrrrr}
\toprule
Period & Avg Diff (bp) & Rel Diff (\%) & Paper Avg Diff (bp) & Paper Rel Diff (\%) \\
\midrule
Original sample (1996--2022) & 8.15 & 3.14 & 7.00 & 2.82 \\
Extended period (2023--2024) & -31.15 & -5.92 & -- & -- \\
Full sample (1996--2024) & 5.44 & 1.95 & -- & -- \\
\bottomrule
\end{tabular}

\caption*{
This table reports summary statistics for the original sample, the extended period,
and the full sample through 2024.
The extension shows that the implied-minus-zero spread turns negative on average
after the end of the original paper's sample.
}
\label{table:figure1_extension_summary}
\end{table}

\section{Conclusion}

We closely replicate Figure 1 from \cite{golez2024holding}
using OptionMetrics end-of-day data together with FRED Treasury rates.
Our results are numerically close to the published values,
with remaining differences likely attributable to differences between
our end-of-day data construction and the paper's intraday CBOE-based implementation.

The extension through 2024 yields a different pattern from that observed in the paper:
the option-implied rate falls below the zero curve rate on average in the post-2022 period.
This result remains after excluding the most extreme negative month in the extension sample,
suggesting that the reversal is not driven entirely by a single outlier.
Overall, the extension indicates that the relationship documented in the original sample
may not be stable in more recent years.

\newpage
\bibliographystyle{jpe}
\bibliography{bibliography.bib}

\end{document}